\documentclass[11pt]{article}
% preamble.tex — shared packages and notation macros (measure theory notes)
\usepackage[T1]{fontenc}
\usepackage{amsmath,amssymb,amsthm,mathtools}
\usepackage[margin=1in]{geometry}
\usepackage[expansion=false]{microtype}
\usepackage{enumitem}
\usepackage[colorlinks=true,linkcolor=black,urlcolor=blue,citecolor=black]{hyperref}

% ---- theorem environments ----
\theoremstyle{plain}
\newtheorem{theorem}{Theorem}[section]
\newtheorem{lemma}[theorem]{Lemma}
\newtheorem{proposition}[theorem]{Proposition}
\newtheorem{corollary}[theorem]{Corollary}
\theoremstyle{definition}
\newtheorem{definition}[theorem]{Definition}
\newtheorem{example}[theorem]{Example}
\theoremstyle{remark}
\newtheorem*{remark}{Remark}

% ---- notation ----
\newcommand{\R}{\mathbb{R}}
\newcommand{\N}{\mathbb{N}}
\newcommand{\salg}{\mathcal{A}}                 % sigma-algebra
\newcommand{\Xset}{X}                            % underlying set
\newcommand{\Prob}{\mathbb{P}}
\newcommand{\E}{\mathbb{E}}
\newcommand{\ind}{\mathbf{1}}                    % indicator
\newcommand{\RN}[2]{\frac{d #1}{d #2}}           % Radon-Nikodym derivative
\newcommand{\abscont}{\ll}                        % absolute continuity
\newcommand{\sing}{\perp}                         % mutual singularity
\newcommand{\given}{\mid}                         % conditioning bar
\newcommand{\dens}{f}                             % density / pmf (cf. likelihood notes)
\newcommand{\restr}{\mathord{\restriction}}      % restriction of a measure
\DeclareMathOperator{\supp}{supp}


\title{Notes on Measure Theory:\\ Absolute Continuity and the Radon--Nikodym Theorem}
\date{\today}

\begin{document}
\maketitle

\begin{abstract}
Results on how one measure sits inside another: absolute continuity and mutual
singularity, the Jordan decomposition of a signed measure, the Radon--Nikodym
theorem with its calculus (linearity, the chain rule, and change of variables),
and the Lebesgue decomposition theorem. The running theme is that a density
$\RN{\nu}{\mu}$ is the measure-theoretic object that both probability densities
and conditional expectations instantiate.
\end{abstract}

\tableofcontents
\bigskip

\section{Absolute continuity and singularity}
\label{sec:abscont}

Throughout, $(\Xset, \salg)$ is a measurable space and all measures are defined
on $\salg$. A measure $\mu$ is \emph{$\sigma$-finite} if $\Xset = \bigcup_n A_n$
with $A_n \in \salg$ and $\mu(A_n) < \infty$; $\sigma$-finiteness is the standing
hypothesis under which the theorems below hold in full generality.

\begin{definition}[Absolute continuity]
\label{def:abscont}
Let $\mu, \nu$ be measures on $(\Xset, \salg)$. We say $\nu$ is
\emph{absolutely continuous} with respect to $\mu$, written $\nu \abscont \mu$,
if
\begin{equation}
  \mu(A) = 0 \;\Longrightarrow\; \nu(A) = 0
  \qquad \text{for every } A \in \salg.
\end{equation}
Equivalently: every $\mu$-null set is $\nu$-null.
\end{definition}

\begin{definition}[Mutual singularity]
\label{def:singular}
The measures $\mu, \nu$ are \emph{mutually singular}, written $\mu \sing \nu$, if
there is a set $E \in \salg$ with
\begin{equation}
  \mu(E) = 0 \qquad \text{and} \qquad \nu(\Xset \setminus E) = 0,
\end{equation}
i.e.\ they live on disjoint sets: $\nu$ is carried by $E$ and $\mu$ by its
complement.
\end{definition}

Absolute continuity and singularity are the two extremes. The Lebesgue
decomposition of Section~\ref{sec:lebesgue} shows every $\sigma$-finite $\nu$
splits uniquely into a piece of each type relative to $\mu$.

\paragraph{The \texorpdfstring{$\varepsilon$--$\delta$}{epsilon-delta} form.}
For \emph{finite} $\nu$, the set-theoretic Definition~\ref{def:abscont} is
equivalent to an analytic continuity statement, which is where the name comes
from.

\begin{proposition}[$\varepsilon$--$\delta$ characterization]
\label{prop:eps-delta}
Let $\nu$ be a finite measure. Then $\nu \abscont \mu$ if and only if for every
$\varepsilon > 0$ there exists $\delta > 0$ such that
\begin{equation}
  A \in \salg,\ \mu(A) < \delta \;\Longrightarrow\; \nu(A) < \varepsilon.
  \label{eq:eps-delta}
\end{equation}
\end{proposition}

\begin{proof}[Proof sketch]
The $\varepsilon$--$\delta$ condition trivially forces $\mu(A)=0 \Rightarrow
\nu(A) = 0$. Conversely, if \eqref{eq:eps-delta} failed there would be
$\varepsilon_0 > 0$ and sets $A_n$ with $\mu(A_n) < 2^{-n}$ but
$\nu(A_n) \ge \varepsilon_0$. Put $A = \limsup_n A_n$. By Borel--Cantelli on
$\mu$, $\mu(A) = 0$, while $\nu(A) \ge \varepsilon_0$ by finiteness of $\nu$ and
reverse Fatou, contradicting $\nu \abscont \mu$.
\end{proof}

\begin{remark}
Finiteness matters in Proposition~\ref{prop:eps-delta}. With
$\mu = $ Lebesgue measure and $\nu(A) = \int_A x^2\,dx$ on $\R$, one has
$\nu \abscont \mu$ (Definition~\ref{def:abscont}) yet no single $\delta$ works
uniformly for the interval family $\{[n, n+\delta]\}$, because $\nu$ is not
finite. The remedy is to test \eqref{eq:eps-delta} on sets of finite $\nu$-measure.
\end{remark}

\paragraph{Signed measures and the Jordan decomposition.}
A \emph{signed measure} $\nu$ assigns values in $(-\infty, +\infty]$ (or
$[-\infty, +\infty)$) and is countably additive. The Hahn decomposition gives a
partition $\Xset = P \sqcup N$ into a positive set $P$ and a negative set $N$;
setting $\nu^+(A) = \nu(A \cap P)$ and $\nu^-(A) = -\nu(A \cap N)$ yields two
genuine (nonnegative) measures.

\begin{proposition}[Jordan decomposition]
\label{prop:jordan}
Every signed measure $\nu$ admits a unique decomposition
\begin{equation}
  \nu = \nu^+ - \nu^-, \qquad \nu^+ \sing \nu^-,
\end{equation}
into mutually singular measures. Its \emph{total variation} is
$\lvert \nu \rvert = \nu^+ + \nu^-$, and $\nu \abscont \mu$ is defined to mean
$\lvert \nu \rvert \abscont \mu$.
\end{proposition}

This reduces every statement about signed measures to the nonnegative case
applied to $\nu^+$ and $\nu^-$ separately, which is how the Radon--Nikodym
theorem is extended to signed and complex measures.

\section{The Radon--Nikodym theorem}
\label{sec:radon-nikodym}

The theorem says that absolute continuity is not merely a null-set condition:
whenever $\nu \abscont \mu$, the measure $\nu$ is obtained from $\mu$ by
integrating a density. It is the converse to the elementary fact that
$A \mapsto \int_A f\,d\mu$ defines a measure absolutely continuous with respect
to $\mu$.

\begin{theorem}[Radon--Nikodym]
\label{thm:rn}
Let $\mu$ and $\nu$ be $\sigma$-finite measures on $(\Xset, \salg)$ with
$\nu \abscont \mu$. Then there exists a measurable function
$f : \Xset \to [0, \infty)$ such that
\begin{equation}
  \nu(A) = \int_A f \, d\mu
  \qquad \text{for all } A \in \salg.
  \label{eq:rn}
\end{equation}
The function $f$ is unique up to $\mu$-a.e.\ equality. It is called the
\emph{Radon--Nikodym derivative} of $\nu$ with respect to $\mu$ and written
\begin{equation}
  f = \RN{\nu}{\mu}.
\end{equation}
\end{theorem}

\begin{remark}[Why $\sigma$-finiteness is needed]
Take $\mu$ the counting measure and $\nu$ the Lebesgue measure on $[0,1]$. Then
$\nu \abscont \mu$ (the only $\mu$-null set is $\emptyset$), yet no $f$
satisfies \eqref{eq:rn}: it would need $\int_{\{x\}} f\,d\mu = f(x) = \nu(\{x\})
= 0$ for every $x$, forcing $f \equiv 0$ and hence $\nu \equiv 0$. Counting
measure is not $\sigma$-finite on an uncountable space, and the theorem fails.
\end{remark}

\begin{remark}[The von Neumann proof]
The slick proof avoids constructing $f$ by hand. Assume $\mu, \nu$ finite, set
$\rho = \mu + \nu$, and consider the bounded linear functional
$g \mapsto \int g \, d\nu$ on the Hilbert space $L^2(\rho)$. By Riesz
representation it equals $\int g h \, d\rho$ for some $h \in L^2(\rho)$; algebra
on \eqref{eq:rn}-type identities shows $0 \le h < 1$ $\mu$-a.e.\ (this uses
$\nu \abscont \mu$), and $f = h / (1 - h)$ is the required derivative. The
$\sigma$-finite case follows by exhausting $\Xset$ with finite pieces.
\end{remark}

\subsection{The calculus of Radon--Nikodym derivatives}

The derivative notation is justified by the following rules, each holding
$\mu$- (respectively $\lambda$-) almost everywhere. They are what make
$\RN{\nu}{\mu}$ behave like an ordinary derivative.

\begin{proposition}[Linearity]
\label{prop:rn-linear}
If $\nu_1, \nu_2 \abscont \mu$ are $\sigma$-finite, then
$\nu_1 + \nu_2 \abscont \mu$ and
\begin{equation}
  \RN{(\nu_1 + \nu_2)}{\mu} = \RN{\nu_1}{\mu} + \RN{\nu_2}{\mu}
  \qquad \mu\text{-a.e.}
\end{equation}
\end{proposition}

\begin{proposition}[Change of variables]
\label{prop:rn-integration}
Let $\nu \abscont \mu$ be $\sigma$-finite with density $f = \RN{\nu}{\mu}$. For
every measurable $g : \Xset \to [0, \infty]$,
\begin{equation}
  \int_\Xset g \, d\nu = \int_\Xset g \, f \, d\mu
  = \int_\Xset g \, \RN{\nu}{\mu} \, d\mu.
  \label{eq:rn-cov}
\end{equation}
The identity extends to $\nu$-integrable $g$ of arbitrary sign.
\end{proposition}

\begin{proof}[Proof idea]
Equation \eqref{eq:rn} is exactly \eqref{eq:rn-cov} for indicators
$g = \ind_A$. Extend by linearity to simple functions, then by monotone
convergence to nonnegative measurable $g$, then split $g = g^+ - g^-$ for the
signed case.
\end{proof}

\begin{proposition}[Chain rule]
\label{prop:rn-chain}
If $\nu \abscont \mu \abscont \lambda$ with all three $\sigma$-finite, then
$\nu \abscont \lambda$ and
\begin{equation}
  \RN{\nu}{\lambda} = \RN{\nu}{\mu} \, \RN{\mu}{\lambda}
  \qquad \lambda\text{-a.e.}
  \label{eq:rn-chain}
\end{equation}
In particular, if $\mu \abscont \nu$ and $\nu \abscont \mu$ (equivalent
measures) then $\RN{\nu}{\mu}$ is strictly positive $\mu$-a.e.\ and
\begin{equation}
  \RN{\mu}{\nu} = \left( \RN{\nu}{\mu} \right)^{-1}
  \qquad \mu\text{-a.e.}
\end{equation}
\end{proposition}

\begin{proof}[Proof idea]
Apply Proposition~\ref{prop:rn-integration} with $g = \ind_A \, \RN{\nu}{\mu}$
and the density $\RN{\mu}{\lambda}$:
$\nu(A) = \int_A \RN{\nu}{\mu}\,d\mu = \int_A \RN{\nu}{\mu}\,\RN{\mu}{\lambda}\,
d\lambda$, so the product is a $\lambda$-density for $\nu$; uniqueness in
Theorem~\ref{thm:rn} gives \eqref{eq:rn-chain}. The reciprocal rule is the case
$\lambda = \nu$, where $\RN{\nu}{\nu} = 1$.
\end{proof}

\section{The Lebesgue decomposition and two instances}
\label{sec:lebesgue}

Radon--Nikodym handles the absolutely continuous case. The Lebesgue
decomposition theorem says that this case, together with the singular one, is
all there is: any $\sigma$-finite measure splits uniquely into an
$\RN{\cdot}{\mu}$-representable part and a part concentrated where $\mu$ is not.

\begin{theorem}[Lebesgue decomposition]
\label{thm:lebesgue}
Let $\mu$ and $\nu$ be $\sigma$-finite measures on $(\Xset, \salg)$. There is a
unique pair of $\sigma$-finite measures $\nu_{\mathrm{ac}}, \nu_{\mathrm{s}}$
with
\begin{equation}
  \nu = \nu_{\mathrm{ac}} + \nu_{\mathrm{s}},
  \qquad \nu_{\mathrm{ac}} \abscont \mu,
  \qquad \nu_{\mathrm{s}} \sing \mu.
  \label{eq:lebesgue}
\end{equation}
Consequently $\nu_{\mathrm{ac}}$ has a density $\RN{\nu_{\mathrm{ac}}}{\mu}$, and
one writes $\RN{\nu}{\mu} := \RN{\nu_{\mathrm{ac}}}{\mu}$ for the
Radon--Nikodym part even when $\nu \not\abscont \mu$.
\end{theorem}

\begin{proof}[Proof idea]
Apply Theorem~\ref{thm:rn} to $\nu$ against the reference $\rho = \mu + \nu$:
$\nu \abscont \rho$, so $\nu = \int h \, d\rho$ for some $0 \le h \le 1$. The set
$E = \{h = 1\}$ carries the singular part $\nu_{\mathrm{s}} = \nu\restr E$, its
complement the absolutely continuous part. Uniqueness follows because a measure
that is both $\abscont \mu$ and $\sing \mu$ must vanish.
\end{proof}

\begin{corollary}
$\nu \abscont \mu$ if and only if $\nu_{\mathrm{s}} = 0$; and $\nu \sing \mu$ if
and only if $\nu_{\mathrm{ac}} = 0$. The two notions of
Section~\ref{sec:abscont} are exactly the vanishing of one summand.
\end{corollary}

\paragraph{Instance 1: probability densities.}
Let $\Prob$ be the law of a real random variable $\Xset$ and let $\mu$ be
Lebesgue measure on $\R$. The distribution decomposes (Theorem~\ref{thm:lebesgue})
into an absolutely continuous part with a probability density function
$p = \RN{\Prob_{\mathrm{ac}}}{\mu}$, an atomic part (a sum of point masses,
singular w.r.t.\ $\mu$), and a singular continuous part (e.g.\ the Cantor
distribution). ``$X$ has density $p$'' is precisely the statement
$\Prob \abscont \mu$ with $\RN{\Prob}{\mu} = p$; then, by
Proposition~\ref{prop:rn-integration},
\begin{equation}
  \E[g(X)] = \int_\R g \, d\Prob = \int_\R g(x)\, p(x)\, dx.
\end{equation}
This is the exact sense in which the likelihood notes' ``density'' is a
Radon--Nikodym derivative: there $\dens(x \given \theta) = \RN{P_\theta}{\mu}$.

\paragraph{Instance 2: conditional expectation.}
Let $(\Omega, \salg, \Prob)$ be a probability space, $\mathcal{G} \subseteq
\salg$ a sub-$\sigma$-algebra, and $X \ge 0$ integrable. Define a measure on
$\mathcal{G}$ by
\begin{equation}
  \nu(G) = \int_G X \, d\Prob, \qquad G \in \mathcal{G}.
\end{equation}
Then $\nu \abscont \Prob\restr\mathcal{G}$ (a $\Prob$-null set in $\mathcal{G}$
is $\nu$-null), so Radon--Nikodym supplies a $\mathcal{G}$-measurable density.

\begin{definition}[Conditional expectation]
The \emph{conditional expectation} of $X$ given $\mathcal{G}$ is
\begin{equation}
  \E[X \mid \mathcal{G}] = \RN{\nu}{\ \Prob\restr\mathcal{G}},
\end{equation}
the Radon--Nikodym derivative on $\mathcal{G}$. It is the unique (up to
$\Prob$-a.s.) $\mathcal{G}$-measurable random variable with
$\int_G \E[X \mid \mathcal{G}] \, d\Prob = \int_G X \, d\Prob$ for all
$G \in \mathcal{G}$.
\end{definition}

Existence and a.s.-uniqueness of conditional expectation are therefore not
separate theorems: they are Theorem~\ref{thm:rn} read on the sub-$\sigma$-algebra
$\mathcal{G}$. The defining averaging identity is the special case $A = G$ of
\eqref{eq:rn}.


\end{document}
